% Declaration of generative AI usage
\section*{Declaration of generative AI and AI-assisted technologies}
During the preparation of this work the author used generative AI
services and tools in order to find resources, research relevant topics,
review prose, and generate code. After using these tools, the author
reviewed and edited the content as needed and takes full responsibility
for the content of this project. In compliance with the code of conduct
for the use of AI at SDU \cite{SDU-AI-CoC}.

\subsection*{Source and literature search}
Claude Code \cite{claude-code} configured with the Claude Opus 4.6 and
Opus 4.7 models \cite{claude-opus} was used as an interactive assistant
to surface candidate references and to cross-check citations against
their sources. An example prompt used during this process looks like
the following:

\begin{lstlisting}[style=prompt]
I am writing the Background chapter of a bachelor thesis on multi-fidelity schedulers (Successive Halving, Hyperband, ASHA) applied to a self-improving coding agent. Can you list the canonical references for each of these methods and verify whether the Karnin et al. 2013 paper is the correct primary citation for Successive Halving?
\end{lstlisting}

\subsection*{Brainstorming and idea generation}
Claude Code \cite{claude-code} was used to discuss design alternatives
for the scheduler implementations and the experimental protocol. Two
recurring use-cases:
\begin{enumerate}
    \item Scheduler design discussions: clarifying the difference
    between synchronous Hyperband and ASHA, sanity-checking the
    promotion-quota formulation, and reviewing the per-rung budget
    accounting. An example prompt used in this stage:
    \begin{lstlisting}[style=prompt]
In ASHA, the promotion quota per rung is computed as max(1, ceil(n / eta^(r+1))). Is this a cumulative cap on total promotions ever from rung r, or a per-rung count refreshed each time a candidate finishes? I want to make sure my implementation matches the original Li et al. 2020 paper.
    \end{lstlisting}
    \item Experimental-protocol brainstorming: identifying threats to
    validity, choosing diagnostic metrics, and deciding which
    constants to fix across runs. An example prompt used in this
    stage:
    \begin{lstlisting}[style=prompt]
I am running one seed per scheduler on swe_verified_mini, four generations each, fixed initial archive. What are the main threats to validity I should disclose, and which diagnostic metrics would help a reader judge whether the cost-vs-accuracy comparison is fair?
    \end{lstlisting}
\end{enumerate}

\subsection*{Programming}
Agentic coding was used during the implementation phase of the project
for editing the fork of the Darwin G\"odel Machine repository, in
particular the scheduler layer and the operational guards described in
Chapter~\ref{chap:implementation}. For this purpose, the tool below has
been used:
\begin{itemize}
    \item Claude Code \cite{claude-code} with the Claude Opus 4.6 and
    Opus 4.7 models \cite{claude-opus} - the only configuration used
    in the project, for planning, implementing, reviewing diffs,
    locating code across the fork, and assisting with debugging of
    Docker, OpenRouter, and SWE-bench harness integration issues.
\end{itemize}
All AI-generated code was reviewed, tested, and integrated by the
author; the author takes full responsibility for the final
implementation.

\subsection*{Language proofreading or linguistic feedback}
Claude Code \cite{claude-code} was used to review and iterate on the
prose of the thesis. The author wrote the chapter body text by hand;
the model was used to flag unclear sentences, inconsistent terminology,
and structural issues in already-written passages, and to suggest
rewordings that the author then accepted, rejected, or rewrote. The
model was not used to ghostwrite chapter body text. An example prompt
used in this stage:

\begin{lstlisting}[style=prompt]
Review this paragraph for clarity and internal consistency with the rest of the Method chapter. Flag any sentence whose meaning is ambiguous, any term used inconsistently with earlier sections, and any claim that is stronger than what the single-seed pilot data supports. Do not rewrite the paragraph; list the issues.
\end{lstlisting}
